\documentclass[11pt,a4paper]{article}

\usepackage[utf8]{inputenc}
\usepackage[T1]{fontenc}
\usepackage[spanish,es-noquoting]{babel}
\usepackage{amsmath,amssymb,amsthm}
\usepackage{mathtools}
\usepackage{booktabs}
\usepackage{graphicx}
\usepackage{natbib}
\usepackage[hidelinks]{hyperref}
\usepackage{geometry}
\geometry{margin=2.6cm}

\theoremstyle{definition}
\newtheorem{definicion}{Definición}[section]
\newtheorem{supuesto}[definicion]{Supuesto}
\newtheorem{requisito}[definicion]{Requisito de medición}
\theoremstyle{plain}
\newtheorem{proposicion}[definicion]{Proposición}
\newtheorem{corolario}[definicion]{Corolario}
\theoremstyle{remark}
\newtheorem{observacion}[definicion]{Observación}

\DeclareMathOperator{\Var}{Var}
\DeclareMathOperator*{\argmax}{arg\,max}
\newcommand{\ind}[1]{\mathbf{1}\!\left[#1\right]}
\newcommand{\solver}{\mathcal{S}}
\newcommand{\Info}{\mathcal{I}}
\newcommand{\clip}{\operatorname{clip}}
\newcommand{\Cl}{\operatorname{Cl}}

\title{\bf Interdependencia epistémica en problemas colaborativos\\
de matemática universitaria:\\
un marco operacional y la validez del cribado agente--agente\\
como sustituto de la interacción estudiante--agente}

\author{%
  Luciano Avegno Cepeda\\
  \small{[afiliación]}\\
  \small{\texttt{lucianoavegno@gmail.com}}
}

\date{Borrador --- \today}

\begin{document}
\maketitle

\begin{abstract}
\noindent
Los problemas de resolución colaborativa (CPS) con partición de información
prometen forzar la co-construcción de conocimiento, pero validar que un problema
efectivamente lo hace exige correr sesiones con estudiantes, un procedimiento caro
que no escala al tamaño de un banco de ejercicios. Proponemos un marco operacional
que reemplaza los indicadores probabilísticos no observables del framework CPP por
estimadores empíricos definidos respecto de un \emph{solver de referencia}
explícito, y demostramos que el índice de balance de carga informativa admite,
en el caso de dos agentes, una forma cerrada equivalente al valor de Shapley que
sólo requiere cuatro proporciones binomiales. Sobre esa base planteamos la
pregunta central: \emph{¿el cribado automático agente--agente predice la calidad
colaborativa observada en la interacción estudiante--agente?} Formulamos un diseño
en dos estudios --- cribado exhaustivo con agentes LLM sobre $50$ problemas y
validación con estudiantes universitarios sobre un submuestreo estratificado de
$25$ --- y especificamos el análisis, la potencia estadística y las amenazas a la
validez antes de la recolección. La contribución esperada es un instrumento de
validación barato para bancos de problemas CPS, junto con evidencia sobre en qué
medida el comportamiento colaborativo de un LLM puede sustituir al de un
estudiante a los fines del diseño de tareas.

\vspace{0.6em}
\noindent\textbf{Palabras clave:} resolución colaborativa de problemas,
interdependencia epistémica, perfil oculto, agentes LLM, validez de constructo,
educación matemática.
\end{abstract}

% =====================================================================
\section{Introducción}
% =====================================================================

Un problema colaborativo con partición de información distribuye entre dos
participantes datos complementarios tales que ninguno alcanza la solución por
separado. La estructura es antigua en psicología social: es el \emph{perfil
oculto} de \citet{stasser1985pooling}, donde el óptimo grupal sólo emerge si los
miembros ponen en común la información no compartida --- y donde el hallazgo
robusto es que, sistemáticamente, no lo hacen. En educación matemática la misma
estructura se usa con signo inverso: no para estudiar el fracaso de la puesta en
común, sino para \emph{forzarla}, obligando a los estudiantes a construir lo que
\citet{roschelle1995construction} llamaron un espacio de problema conjunto.

El marco de referencia para evaluar esta competencia es el de PISA 2015
\citep{oecd2017cps}, que cruza cuatro procesos individuales de resolución
(explorar y comprender; representar y formular; planificar y ejecutar; monitorear
y reflexionar) con tres competencias colaborativas (establecer y mantener
comprensión compartida; tomar la acción apropiada; establecer y mantener la
organización del equipo), produciendo una matriz de $12$ celdas.

El problema práctico es de escala. Diseñar un banco de problemas colaborativos
requiere verificar, para cada ítem, que la partición sea genuina: que el dato $A$
solo no determine la respuesta, que el dato $B$ solo tampoco, y que la resolución
conjunta exija más de una transacción trivial. Verificarlo con estudiantes es
caro, lento y consume la novedad del ítem. Verificarlo por juicio experto es
barato pero subjetivo y no auditable.

Este trabajo explora una tercera vía: usar dos agentes LLM con la misma partición
de información como \emph{instrumento de cribado}, computar sobre sus transcripts
una batería de indicadores, y preguntar si ese cribado predice lo que ocurre
cuando el segundo agente es reemplazado por un estudiante universitario.

La pregunta tiene un antecedente directo. \citet{herborn2020agents} analizaron si
los agentes computacionales de PISA 2015 capturan válidamente la dinámica de la
interacción humana, y no hallaron diferencias significativas en desempeño según el
tipo de compañero, aunque sí más acciones cuando el compañero era un par humano.
Su estudio, sin embargo, compara \emph{humano--agente} contra
\emph{humano--humano} con agentes guionados por reglas. Nosotros comparamos
\emph{agente--agente} contra \emph{humano--agente}, con agentes LLM
no guionados, y con un objetivo distinto: no evaluar al estudiante, sino
\textbf{validar el problema}. En el mismo sentido, \citet{peng2025asymmetry}
muestran que agentes LLM bajo asimetría de información mejoran sustancialmente con
protocolos explícitos de comunicación y verificación, lo que sugiere que su
comportamiento colaborativo por defecto dista de ser óptimo --- exactamente la
variabilidad que un instrumento de cribado necesita para discriminar problemas.

\paragraph{Contribuciones.}
\begin{enumerate}
  \item Una reformulación operacional del framework CPP (§\ref{sec:marco}) que
        sustituye probabilidades y entropías no observables por estimadores
        binomiales relativos a un solver de referencia declarado, con propiedades
        de acotación demostradas y varianzas explícitas.
  \item Una forma cerrada del balance de carga informativa
        (Proposición~\ref{prop:ibc}) que elimina la entropía condicional del
        marco y reduce su cómputo a cuatro proporciones.
  \item La identificación de dos defectos formales en el índice compuesto de
        rendimiento colaborativo y su reemplazo por un reporte vectorial
        (§\ref{sec:cy}).
  \item Un diseño experimental preregistrable en dos estudios, con hipótesis,
        potencia y análisis especificados de antemano (§\ref{sec:diseno}).
  \item Una plataforma de código abierto que instrumenta el protocolo
        (§\ref{sec:implementacion}).
\end{enumerate}

% =====================================================================
\section{Marco formal}\label{sec:marco}
% =====================================================================

\subsection{Problemas con partición de información}

\begin{definicion}[Instancia CPS]\label{def:instancia}
Una \emph{instancia CPS con partición} es una tupla
$\mathcal{P}=(Q, d_A, d_B, a^\ast, \mathcal{A})$ donde $Q$ es el enunciado público,
$d_A$ y $d_B$ son los datos privados asignados a cada participante,
$\mathcal{A}$ es el espacio de respuestas admisibles y $a^\ast\in\mathcal{A}$ es
la respuesta correcta única. Para $S\subseteq D:=\{d_A,d_B\}$ escribimos
$\Info(S) := Q\cup S$ para el estado de información correspondiente.
\end{definicion}

La unicidad de $a^\ast$ no es un tecnicismo: sin ella, la corrección no es
verificable automáticamente y todo el aparato de medición que sigue se vuelve
dependiente de juicio humano.

\subsection{Solver de referencia y función de competencia}

El defecto central del framework CPP original es que su indicador principal se
apoya en $P(\text{solución}\mid \Info)$, una cantidad que no está definida sobre
ningún espacio de probabilidad especificado. Para un problema de matemática con
respuesta única, «la probabilidad de la solución dada cierta información» no
existe hasta que se fija \emph{quién} intenta resolverlo. Lo hacemos explícito.

\begin{definicion}[Solver de referencia]\label{def:solver}
Un \emph{solver de referencia} es un mapa estocástico
$\solver:\Info\to\mathcal{A}$, especificado por completo (modelo, versión,
temperatura, prompt, presupuesto de razonamiento). Su \emph{función de
competencia} es el juego cooperativo $(D,v)$ con
\[
  v(S) \;:=\; \Pr\big[\solver(\Info(S)) = a^\ast\big],
  \qquad S\subseteq D, \qquad v:2^{D}\to[0,1].
\]
\end{definicion}

\begin{definicion}[Estimador]\label{def:estimador}
Dadas $N$ ejecuciones independientes de $\solver$ sobre $\Info(S)$,
\[
  \hat v(S) \;=\; \frac{1}{N}\sum_{j=1}^{N}\ind{\solver_j(\Info(S)) = a^\ast},
  \qquad N\hat v(S)\sim \mathrm{Bin}\big(N, v(S)\big),
\]
con intervalo de confianza de Wilson \citep{wilson1927probable}. La corrección se
determina por equivalencia simbólica contra $a^\ast$, no por autorreporte.
\end{definicion}

\begin{observacion}[Relatividad al solver, y por qué es aceptable]
$v$ depende de $\solver$. Esto no es un defecto sino una explicitación: toda
medida de dificultad es relativa a una población de resolutores, igual que la
dificultad de un ítem en teoría de respuesta al ítem es relativa a la muestra de
calibración. Lo que exigimos es que $\solver$ se declare y que la estabilidad del
\emph{ordenamiento} de problemas entre solvers distintos se reporte
(§\ref{sec:amenazas}).
\end{observacion}

\begin{supuesto}[Monotonía informacional]\label{sup:monotonia}
$v(\emptyset)\le v(\{d_i\})\le v(D)$ para $i\in\{A,B\}$.
\end{supuesto}

El Supuesto~\ref{sup:monotonia} es normativo, no empírico: un resolutor ideal no
empeora al recibir información adicional. Los LLM pueden violarlo por efectos de
distracción. Las violaciones no se descartan, se reportan: son un diagnóstico
sobre el ítem o sobre el solver.

\subsection{Interdependencia epistémica}

\begin{definicion}[Ganancia alcanzable y no degeneración]
$\Delta := v(D)-v(\emptyset)$. La instancia es \emph{no degenerada} si
$\Delta>0$.
\end{definicion}

\begin{definicion}[Interdependencia epistémica de diseño]\label{def:ie0}
\[
  IE_0 \;:=\; 1-\frac{\max_i\big[v(\{d_i\})-v(\emptyset)\big]}{\Delta}.
\]
\end{definicion}

\begin{proposicion}\label{prop:ie}
Bajo el Supuesto~\ref{sup:monotonia} y no degeneración:
\emph{(i)} $IE_0\in[0,1]$;
\emph{(ii)} $IE_0=1 \iff v(\{d_A\})=v(\{d_B\})=v(\emptyset)$;
\emph{(iii)} $IE_0=0 \iff \max_i v(\{d_i\})=v(D)$.
\end{proposicion}

\begin{proof}
Por monotonía, $0\le v(\{d_i\})-v(\emptyset)\le \Delta$ para cada $i$, luego el
cociente pertenece a $[0,1]$ y (i) sigue. (ii) y (iii) son los casos de igualdad
en cada extremo del cociente.
\end{proof}

\begin{observacion}[Por qué la normalización no es cosmética]\label{obs:normalizacion}
El indicador original $IE = 1-\max_i P(\text{sol}\mid I_i)$ confunde
\emph{dificultad} con \emph{interdependencia}. Un problema que nadie resuelve ni
siquiera con ambos datos ($v(D)=v(\{d_i\})=0$) alcanza el máximo del indicador sin
ser colaborativo en absoluto: es simplemente imposible. Condicionar sobre la
ganancia alcanzable $\Delta$ separa ambos constructos, y deja la dificultad como
un indicador aparte ($1-v(D)$) en vez de contaminar al principal.
\end{observacion}

\begin{definicion}[Interdependencia dinámica y persistencia]
Sea $m_{1:t}$ el transcript público hasta la ronda $t$ e
$I_{i,t} := Q\cup\{d_i\}\cup m_{1:t}$. Con
$v_t^{(i)} := \Pr[\solver(I_{i,t})=a^\ast]$,
\[
  IE_t \;:=\; \clip_{[0,1]}\!\left(1-\frac{\max_i v_t^{(i)}-v(\emptyset)}{\Delta}\right),
  \qquad
  t_{\mathrm{crit}} \;:=\; \max\{\,t\le T : IE_t\ge \tau\,\},
\]
con $\tau=1/2$ y la convención $\max\emptyset=0$.
\end{definicion}

Un problema cuya partición se disuelve en el primer intercambio tiene
$t_{\mathrm{crit}}=0$. La condición prescriptiva de diseño es
$t_{\mathrm{crit}}\ge 2$.

\subsection{Balance de carga informativa}

Un problema puede exigir que ambos hablen y aun así ser colaborativamente
superficial, si uno aporta el $95\%$ del peso deductivo y el otro un dato trivial.
El framework original captura esto con entropías condicionales, lo que introduce
dos problemas: la entropía diferencial sobre un espacio de respuestas continuo es
dependiente de escala y puede ser negativa, con lo cual el índice no queda acotado
en $[0,1]$; y en ningún caso es estimable con el mismo experimento que $IE$.
Reemplazamos la entropía por el valor de Shapley \citep{shapley1953value}, en la
misma línea en que se lo usa para atribución en aprendizaje automático
\citep{lundberg2017unified}.

\begin{definicion}[Contribución marginal]
Para el juego $(D,v)$ con $|D|=2$,
\[
  \varphi_A \;=\; \tfrac12\big[v(\{d_A\})-v(\emptyset)\big] \;+\;
                  \tfrac12\big[v(D)-v(\{d_B\})\big],
\]
y simétricamente $\varphi_B$.
\end{definicion}

\begin{proposicion}\label{prop:ibc}
Se cumple que
\emph{(i)} $\varphi_A+\varphi_B=\Delta$ \emph{(eficiencia)}; y
\emph{(ii)} $\varphi_A-\varphi_B = v(\{d_A\})-v(\{d_B\})$.
\end{proposicion}

\begin{proof}
Para (i), sumando las dos expresiones,
$\varphi_A+\varphi_B=\tfrac12[v(\{d_A\})+v(\{d_B\})-2v(\emptyset)]
+\tfrac12[2v(D)-v(\{d_A\})-v(\{d_B\})]=v(D)-v(\emptyset)=\Delta$.
Para (ii), restando,
\[
\varphi_A-\varphi_B=\tfrac12\big[v(\{d_A\})-v(\{d_B\})\big]
+\tfrac12\big[v(\{d_A\})-v(\{d_B\})\big]=v(\{d_A\})-v(\{d_B\}). \qedhere
\]
\end{proof}

\begin{definicion}[Índice de balance de carga informativa]
\[
  IBC \;:=\; 1-\frac{|\varphi_A-\varphi_B|}{\Delta}
        \;=\; 1-\frac{\big|v(\{d_A\})-v(\{d_B\})\big|}{v(D)-v(\emptyset)}.
\]
\end{definicion}

\begin{corolario}
Bajo el Supuesto~\ref{sup:monotonia} y no degeneración, $IBC\in[0,1]$, con
$IBC=1$ si y sólo si ambos datos producen idéntica ganancia individual, e
$IBC=0$ si y sólo si uno de los datos concentra toda la ganancia alcanzable.
\end{corolario}

\begin{proof}
Sin pérdida de generalidad $v(\{d_A\})\ge v(\{d_B\})$. Por monotonía
$v(\{d_A\})\le v(D)$ y $v(\{d_B\})\ge v(\emptyset)$, luego
$0\le v(\{d_A\})-v(\{d_B\})\le v(D)-v(\emptyset)=\Delta$. Los casos de igualdad
dan las dos caracterizaciones.
\end{proof}

La Proposición~\ref{prop:ibc}(ii) es el resultado operativamente relevante: el
$IBC$ se computa con las mismas cuatro proporciones que ya se estimaron para
$IE_0$, sin ningún experimento adicional y sin entropías.

\begin{observacion}[Extensión a $k>2$ agentes]
Con $k$ participantes, $\varphi_i$ recupera la forma general de Shapley y el
índice se generaliza como $IBC_k = 1-(\max_i\varphi_i-\min_i\varphi_i)/\Delta$,
o mediante el coeficiente de Gini del vector $(\varphi_1,\dots,\varphi_k)$
si interesa penalizar desbalances distribuidos y no sólo el rango.
\end{observacion}

\subsection{Incertidumbre de los estimadores}

\begin{proposicion}[Varianza por método delta]\label{prop:delta}
Sea $g(v)=\big(v_A-v_B\big)/\big(v_D-v_\emptyset\big)$ con estimaciones
independientes de tamaño $N$ cada una. Entonces
\[
  \Var(\hat g)\;\approx\;
  \frac{v_A(1-v_A)+v_B(1-v_B)}{N\,\Delta^{2}}
  \;+\;
  \frac{(v_A-v_B)^{2}\,\big[v_D(1-v_D)+v_\emptyset(1-v_\emptyset)\big]}{N\,\Delta^{4}}.
\]
\end{proposicion}

\begin{proof}
Las derivadas parciales son $\partial g/\partial v_A = 1/\Delta$,
$\partial g/\partial v_B = -1/\Delta$,
$\partial g/\partial v_D = -(v_A-v_B)/\Delta^{2}$ y
$\partial g/\partial v_\emptyset = (v_A-v_B)/\Delta^{2}$. Aplicando el método
delta con $\Var(\hat v_S)=v_S(1-v_S)/N$ e independencia entre celdas se obtiene la
expresión.
\end{proof}

El término en $\Delta^{-4}$ advierte que la estimación se vuelve inestable cuando
la ganancia alcanzable es pequeña. Cuando $\hat\Delta$ es chico recomendamos
intervalos bootstrap BCa en lugar de la aproximación normal.

\paragraph{Tamaño muestral.} Verificar el Supuesto de insuficiencia individual
($v(\{d_i\})\le\delta$) es un problema de cota superior con cero eventos. Por la
regla de tres, si ninguna de $N$ ejecuciones acierta, la cota superior unilateral
al $95\%$ es aproximadamente $3/N$. Así, $N=30$ certifica $v(\{d_i\})\lesssim0.10$
y $N=60$ certifica $\lesssim 0.05$. Fijamos $N=30$ por celda como valor por
defecto, elevándolo a $60$ en los ítems que se reporten como limítrofes.

\subsection{Estructura de la solución: emergencia y cota de rondas}

\begin{definicion}[DAG de solución]
La resolución canónica de $\mathcal{P}$ se representa como un DAG $G=(V,E)$ cuyos
nodos son proposiciones o pasos deductivos y cuyas aristas son dependencias.
Para $S\subseteq D$, la \emph{clausura} $\Cl(S)\subseteq V$ es el conjunto de
nodos derivables a partir de $\Info(S)$ recorriendo $E$.
\end{definicion}

\begin{definicion}[Métrica de emergencia epistémica]
\[
  MEE \;:=\; \frac{\big|\,V\setminus\big(\Cl(\{d_A\})\cup \Cl(\{d_B\})\big)\big|}
                  {\big|\,V\setminus \Cl(\emptyset)\big|}\;\in[0,1].
\]
\end{definicion}

$MEE$ mide la fracción de pasos no triviales que no pertenecen al dominio de
ninguno de los dos agentes por separado, es decir, que sólo existen como producto
de la deducción conjunta. La normalización por $V\setminus\Cl(\emptyset)$ excluye
los pasos derivables del enunciado público, que no informan sobre la partición.

\begin{definicion}[Cota inferior de rondas]\label{def:tmin}
Asígnese a cada nodo $\nu\in V$ el lado $\lambda(\nu)\in\{A,B,AB\}$ del dato
privado mínimo que su derivación requiere. Definimos
\[
  t_{\min} \;:=\; \max_{\pi\,\in\,\mathrm{caminos}(G)}
  \#\big\{\, j : \lambda(\pi_j)\neq\lambda(\pi_{j+1}),\;
  \lambda(\pi_j),\lambda(\pi_{j+1})\in\{A,B\} \,\big\},
\]
el máximo número de alternancias de lado a lo largo de un camino dirigido de $G$.
\end{definicion}

$t_{\min}$ es una cota inferior estructural sobre la cantidad de intercambios
necesarios: cada alternancia de lado en la cadena de dependencias exige al menos
una transmisión de información. Es un dato de \emph{diseño}, computable del DAG,
y es lo que vuelve interpretable la eficiencia observada.

\begin{definicion}[Eficiencia del episodio]
$E := t_{\min}/t_{\mathrm{obs}} \in(0,1]$, con $t_{\mathrm{obs}}$ el número de
rondas efectivamente usadas hasta la primera respuesta correcta.
\end{definicion}

\subsection{Divulgación prematura}

El colapso del split en el primer turno --- un agente vuelca todos sus datos
privados de entrada --- destruye la fase exploratoria. Sostenemos que debe
\textbf{medirse}, no prohibirse por prompt: una restricción explícita en el prompt
sustituye el fenómeno de interés (el comportamiento comunicativo espontáneo del
modelo) por otro (su obediencia a la restricción).

\begin{definicion}[Tiempo de divulgación]
$T_i := \min\{t : d_i \text{ es recuperable en su totalidad de } m_{1:t}\}$,
determinado por un clasificador de implicación auditado contra codificación
humana. El \emph{índice de divulgación prematura} es $PD_i := \ind{T_i=1}$ y la
\emph{tasa de fuga} de una condición es $\mathbb{E}[PD_i]$ sobre sus episodios.
\end{definicion}

Consecuentemente, la restricción de no divulgar en el turno~1 se incorpora como
\emph{condición experimental} cruzada, no como parche del instrumento.

\subsection{Calidad del proceso y el juez automático}

\begin{definicion}[CQI]
Sobre las $12$ celdas de PISA \citep{oecd2017cps}, con
$q_k\in\{0,1,2,3\}$ en escala ordinal desde ausencia de interacción hasta
emergencia de entendimiento conjunto,
$CQI := \tfrac{1}{36}\sum_{k=1}^{12} q_k \in[0,1]$.
\end{definicion}

La codificación manual no escala a centenares de transcripts. Adoptamos un juez
LLM, cuya validez no se presume sino que se estima. \citet{zheng2023judging}
documentan que jueces de frontera alcanzan acuerdo superior al $80\%$ con
evaluadores humanos, comparable al acuerdo humano--humano, pero también sesgos de
posición, de verbosidad y de autopreferencia. El protocolo mínimo que adoptamos:

\begin{requisito}[Protocolo de juez]\label{req:juez}
\emph{(a)} Al menos dos jueces de familias de modelos distintas, ninguno de la
familia que generó los transcripts;
\emph{(b)} ciego a la condición experimental y al identificador del problema;
\emph{(c)} orden de presentación aleatorizado;
\emph{(d)} acuerdo juez--humano reportado como $\alpha$ de Krippendorff ordinal
\citep{krippendorff2018content} sobre una submuestra estratificada del $20\%$
codificada por dos anotadores humanos entrenados. Se adoptan los umbrales
convencionales $\alpha\ge0.80$ para conclusiones firmes y
$0.667\le\alpha<0.80$ para conclusiones tentativas.
\end{requisito}

\begin{requisito}[Corrección verificada]\label{req:correccion}
$\ind{\text{correcto}}$ se determina por equivalencia simbólica de la respuesta
final contra $a^\ast$ mediante un sistema de álgebra computacional, con
normalización previa. En ningún caso se acepta la autodeclaración de resolución
por parte del agente o del participante.
\end{requisito}

El Requisito~\ref{req:correccion} no es pedantería: un marcador de resolución
autodeclarado es falsificable por el propio sujeto medido, y convierte la variable
dependiente principal en una medida de disposición a declararse exitoso.

\subsection{Crítica del rendimiento colaborativo compuesto}\label{sec:cy}

El framework original define
$CY = 0.35\,CDI + 0.45\,\ind{\text{correcto}} + 0.20\,\kappa$. Señalamos dos
defectos formales.

\begin{observacion}[Varianza intra-problema nula]
$CDI$ es, por su propia definición, \emph{prescriptivo}: una función de
$\mathcal{P}$ y no del episodio $\omega$. Por lo tanto
$\Var_\omega(CY\mid\mathcal{P}) = \Var_\omega\big(0.45\,\ind{\text{correcto}}
+0.20\,\kappa \mid \mathcal{P}\big)$: el término $0.35\,CDI$ no aporta varianza
alguna dentro de un problema y opera únicamente como desplazamiento entre
problemas. Un índice de rendimiento a nivel de episodio no puede contener un
término constante en el episodio. La corrección mínima es sustituir $CDI$ por su
contraparte descriptiva $CQI$.
\end{observacion}

\begin{observacion}[Escalarización con pérdida]
El propio framework afirma que calidad de proceso y exactitud del resultado son
empíricamente ortogonales. Si lo son, una suma ponderada de peso fijo colapsa dos
constructos independientes en un escalar no identificable: infinitos pares
$(CQI,\ind{\text{correcto}})$ producen el mismo $CY$, y la ortogonalidad garantiza
que esa pérdida no es recuperable. Adicionalmente, $\kappa$ no está definido en el
documento base.
\end{observacion}

\begin{definicion}[Reporte vectorial]\label{def:vectores}
En lugar de un compuesto, reportamos dos perfiles:
\[
  \Phi(\mathcal{P}) = \big(IE_0,\; t_{\mathrm{crit}},\; IBC,\; MEE,\; CDI,\; 1-v(D)\big),
  \qquad
  \Psi(\omega) = \big(CQI,\; \ind{\text{correcto}},\; E,\; PD\big),
\]
de diseño y de episodio respectivamente. Todo índice compuesto se reporta como
secundario y acompañado de análisis de sensibilidad sobre sus pesos.
\end{definicion}

\begin{observacion}[Redundancia entre indicadores]
Cinco indicadores correlacionados no son cinco constructos. Reportamos la matriz
de correlaciones de $\Phi$ y un análisis de componentes principales; si la primera
componente explica más del $80\%$ de la varianza, el marco debe reducirse.
\end{observacion}

% =====================================================================
\section{Diseño experimental}\label{sec:diseno}
% =====================================================================

\subsection{Pregunta y hipótesis}

\begin{quote}
\emph{¿En qué medida el cribado automático agente--agente de un problema
colaborativo predice la calidad de la colaboración observada cuando el segundo
participante es un estudiante universitario?}
\end{quote}

\begin{description}
  \item[H1 (transferencia de ordenamiento).] $\rho_s\big(CQI^{AA}(p),\,CQI^{HA}(p)\big)>0$
        entre problemas, con $\rho_s$ el coeficiente de Spearman.
  \item[H2 (validez predictiva).] $\Phi^{AA}$ predice $CQI^{HA}$ con
        $R^2$ validado cruzadamente mayor que cero.
  \item[H3 (utilidad de decisión).] El cribado, usado como clasificador binario
        de «problema colaborativamente adecuado», alcanza $AUC>0.70$ contra la
        etiqueta derivada del estudio humano.
  \item[H4 (divergencia comportamental).] La tasa de fuga $\mathbb{E}[PD]$ es
        mayor en agente--agente que en estudiante--agente. Esta hipótesis es
        direccional y su confirmación acota el alcance de H1.
  \item[H5 (validez convergente).] La rúbrica experta a priori de cinco
        dimensiones correlaciona positivamente con los indicadores computados de
        $\Phi$.
\end{description}

H4 merece énfasis: es la hipótesis que puede falsar el programa. Si los LLM
colapsan el split sistemáticamente y los estudiantes no, el cribado mide un
fenómeno distinto del que pretende predecir, y H1 debería fallar.

\subsection{Estudio 1: cribado agente--agente}

\textbf{Materiales.} $50$ problemas con partición, distribuidos sobre seis áreas
de matemática universitaria (álgebra lineal, introducción al álgebra, introducción
al cálculo, cálculo diferencial e integral, cálculo en varias variables,
ecuaciones diferenciales ordinarias), cada uno con su $a^\ast$ estructurada y su
DAG anotado.

\textbf{Diseño.} Factorial $2\times2$ intra-ítem: familia de modelo
$\{M_1,M_2\}$ $\times$ restricción de divulgación $\{$libre, restringida$\}$, con
$R=8$ réplicas por celda. Total: $50\times4\times8=1600$ episodios.

\textbf{Medición.} Para cada problema, $\Phi$ vía $N=30$ ejecuciones del solver de
referencia sobre cada una de las cuatro celdas de información. Para cada episodio,
$\Psi$ vía el protocolo del Requisito~\ref{req:juez}.

\subsection{Estudio 2: validación estudiante--agente}

\textbf{Muestra de ítems.} Submuestreo estratificado de $25$ problemas por
terciles de $IE_0$ y $t_{\mathrm{crit}}$, garantizando cobertura del rango del
cribado y no sólo de su extremo favorable.

\textbf{Participantes.} Estudiantes de grado con la materia correspondiente
aprobada o en curso. Cada participante resuelve $3$ problemas asignados
aleatoriamente, en modo estudiante--agente, con el agente en el rol del
poseedor de $d_B$.

\textbf{Tamaño.} Con $n_s$ estudiantes y $3$ episodios cada uno, el número de
episodios por problema es $3n_s/25$. La confiabilidad de la media por problema
sigue Spearman--Brown, $\rho_k = k\,\mathrm{ICC}/[1+(k-1)\mathrm{ICC}]$. Con
$\mathrm{ICC}=0.30$ y $n_s=60$ resulta $k=7.2$ y $\rho_k\approx0.76$, adecuado
para correlacionar medias por problema.

\subsection{Análisis planificado}

El modelo principal del Estudio 2 es de efectos aleatorios cruzados,
\[
  CQI_{sp} \;=\; \beta_0 + \boldsymbol{\beta}^{\!\top}\Phi_p + u_s + w_p + \varepsilon_{sp},
  \qquad u_s\sim\mathcal{N}(0,\sigma_s^2),\;\; w_p\sim\mathcal{N}(0,\sigma_p^2),
\]
con estudiantes y problemas como factores aleatorios cruzados.

\paragraph{Potencia para H1.} Por la transformación $z$ de Fisher, el número de
problemas necesario es
$n = \big[(z_{1-\alpha/2}+z_{1-\beta})/\operatorname{artanh}\rho\big]^2+3$.
Para $\rho=0.5$, $\alpha=0.05$ bilateral y potencia $0.80$: $n\approx29$. Para
$\rho=0.4$: $n\approx47$. El Estudio 1 se dimensiona en $50$ problemas para
detectar correlaciones moderadas; el Estudio 2 en $25$ conserva potencia sólo para
$\rho\gtrsim0.53$, límite que se declara explícitamente.

\paragraph{Corrección por atenuación.} Ambos términos de H1 son medidas con
error. Reportamos la correlación observada y su versión desatenuada
\citep{spearman1904proof},
$\rho_{\text{true}} = \rho_{\text{obs}}/\sqrt{\alpha_{AA}\,\alpha_{HA}}$,
con $\alpha_\bullet$ las confiabilidades estimadas de cada medida. Omitir esta
corrección subestima sistemáticamente la transferencia.

% =====================================================================
\section{Amenazas a la validez}\label{sec:amenazas}
% =====================================================================

\paragraph{Contaminación del corpus de entrenamiento.} Si un problema aparece en
los datos de preentrenamiento, $v(\{d_i\})$ se infla y $IE_0$ se subestima.
Mitigaciones: ítems de autoría propia; prueba canario consistente en perturbar
constantes numéricas y verificar que $\hat v$ se mueva como predice la estructura;
reporte de $\hat v(\emptyset)$, cuyo exceso sobre el azar es señal directa de
contaminación o de un enunciado que filtra la respuesta.

\paragraph{Relatividad al solver.} Todo $\Phi$ depende de $\solver$. Se replica el
cribado con al menos dos solvers de familias distintas y se reporta la
concordancia de ordenamientos mediante $\tau$ de Kendall. Si $\tau$ es baja, los
indicadores no son propiedades del problema sino del modelo, y el programa debe
reformularse.

\paragraph{Sesgos del juez.} Cubiertos por el Requisito~\ref{req:juez}. Se reporta
adicionalmente la correlación entre longitud del transcript y $CQI$ asignado como
prueba directa de sesgo de verbosidad.

\paragraph{Validez de constructo del LLM como sustituto de estudiante.} Es
precisamente lo que el estudio pone a prueba y no lo que asume. Un resultado nulo
en H1 es informativo y publicable.

\paragraph{Efectos de aprendizaje y orden.} Contrabalanceo del orden de los tres
problemas por participante mediante cuadrado latino.

\paragraph{Ética.} Consentimiento informado; participación desacoplada de toda
calificación; seudonimización en el momento de la captura; derecho de retiro con
borrado; retención y acceso a los datos declarados por anticipado. Sometido a
comité de ética institucional antes del Estudio 2.

% =====================================================================
\section{Implementación}\label{sec:implementacion}
% =====================================================================

El protocolo se instrumenta sobre una plataforma web propia (FastAPI y React) que
ya soporta los tres modos de interacción --- agente--agente, estudiante--agente y
acceso a la resolución --- y persiste transcripts completos con rol y orden. La
Tabla~\ref{tab:brecha} enumera lo que el protocolo exige y el estado actual de la
implementación.

\begin{table}[t]
\centering
\small
\begin{tabular}{@{}lll@{}}
\toprule
\textbf{Constructo} & \textbf{Requiere} & \textbf{Estado} \\
\midrule
Transcripts por rol y orden        & tabla de mensajes                  & implementado \\
Modos AA y HA                      & dos modos de sesión                & implementado \\
Partición $d_A/d_B$                & campos privados por problema       & implementado \\
$\ind{\text{correcto}}$ (Req.~\ref{req:correccion}) & $a^\ast$ estructurada + verificador simbólico & \textbf{pendiente} \\
$v(S)$, $IE_t$, $IBC$              & arnés de solver de referencia      & \textbf{pendiente} \\
$MEE$, $t_{\min}$                  & DAG de solución anotado            & \textbf{pendiente} \\
$CQI$                              & pipeline de juez y vector $q$      & \textbf{pendiente} \\
Condiciones y asignación           & aleatorización registrada          & \textbf{pendiente} \\
$E$, costos                        & tokens y latencia por turno        & \textbf{pendiente} \\
Reproducibilidad                   & versionado de ítems y prompts      & \textbf{pendiente} \\
Análisis                           & exportación JSONL                  & \textbf{pendiente} \\
Ética                              & consentimiento y seudonimización   & \textbf{pendiente} \\
\bottomrule
\end{tabular}
\caption{Correspondencia entre constructos del marco y capacidades de la
plataforma.}
\label{tab:brecha}
\end{table}

% =====================================================================
\section{Discusión}
% =====================================================================

El aporte que este diseño puede hacer es asimétrico en el buen sentido: ambos
resultados posibles son informativos. Si el cribado agente--agente predice la
calidad colaborativa humano--agente, se obtiene un instrumento barato para validar
bancos de problemas CPS a escala, con criterios explícitos y auditables en lugar
de juicio experto. Si no la predice, se obtiene evidencia cuantitativa de que el
comportamiento colaborativo de los LLM diverge del humano en una dimensión
concreta y medible --- plausiblemente la divulgación prematura de información
privada, que H4 aísla --- lo que acota el uso de agentes sintéticos como
sustitutos de estudiantes en investigación educativa.

Por fuera de ese resultado, la reformulación del marco tiene valor independiente:
al anclar los indicadores en un solver de referencia declarado, los vuelve
replicables por terceros, algo que las definiciones probabilísticas originales no
permitían.

\bibliographystyle{apalike}
\bibliography{referencias}

\end{document}
